\documentclass{wscpaperproc}
\usepackage{latexsym}      % extra symbols (WSC template standard)
\usepackage{mathptmx}      % Times New Roman for math (WSC template standard)
\usepackage[T1]{fontenc}   % proper font encoding (WSC template standard)
\usepackage{booktabs}      % professional table rules
\usepackage{listings}      % code listings
\usepackage{microtype}     % better text justification
\usepackage{url}           % URL formatting
\usepackage{amsmath}       % math
\usepackage{graphicx}      % figures
\usepackage{pifont}        % \cmark / \xmark
\usepackage{multirow}      % table cell merging
\usepackage[pdftex,colorlinks=true,urlcolor=blue,citecolor=black,anchorcolor=black,linkcolor=black]{hyperref}

\graphicspath{{figs/}}

% ---- Code listing style ---
\lstset{
  basicstyle=\ttfamily\small,
  columns=flexible,
  keepspaces=true,
  frame=single,
  framerule=0.4pt,
  xleftmargin=6pt,
  xrightmargin=6pt,
  breaklines=true,
  showstringspaces=false,
  numbers=none
}

\newcommand{\cmark}{\ding{51}}
\newcommand{\xmark}{\ding{55}}

% ============================================================
\begin{document}

\WSCpagesetup{ZHANG, NIE, and ZHAO}

\title{HEAS: HIERARCHICAL EVOLUTIONARY AGENT-BASED SIMULATION FRAMEWORK
       FOR MULTI-OBJECTIVE POLICY SEARCH}

\author{\begin{center}Ruiyu ZHANG\textsuperscript{1}, Lin NIE\textsuperscript{2}, and Xin ZHAO\textsuperscript{2}\\
[11pt]
\textsuperscript{1}Dept.~of Politics and Public Administration, The University of Hong Kong, Hong Kong SAR, CHINA\\
\textsuperscript{2}Dept.~of Applied Social Sciences, The Hong Kong Polytechnic University, Hong Kong SAR, CHINA\end{center}
}

\maketitle

\vspace{-12pt}

% ============================================================
\begin{abstract}
Metric aggregation divergence is a hidden confound in agent-based model policy
search: when optimization, tournament evaluation, and statistical validation
independently implement outcome metric extraction, champion selection reflects
aggregation artifact rather than policy quality. We propose Hierarchical Evolutionary Agent Simulation (HEAS), a composable
framework that eliminates this confound through a runtime-enforceable metric
contract---a uniform metrics\_episode() callable shared identically by all
pipeline stages. Removing the confound yields robust champion selection: in a controlled experiment (n=30), HEAS reduces rank reversals by 50\%
relative to ad-hoc aggregation (Cohen h=0.215); the HEAS champion wins all 32
held-out ecological scenarios---a null-safety result that would be
uninterpretable under aggregation divergence. The contract additionally reduces
coupling code by 97\% (160 to 5 lines) relative to Mesa~3.3.1. Three case
studies validate composability across ecological, enterprise, and mean-field ordinary differential equation
dynamics.
\end{abstract}

% ============================================================
\section{INTRODUCTION}

Agent-based models are foundational for policy design across ecology,
economics, and organizational science, yet coupling an ordinary differential equation-style or tabular ABM
to evolutionary optimization requires integration infrastructure that
general-purpose frameworks do not provide. This paper addresses a specific
integration failure: when researchers link an ABM to an evolutionary algorithm,
the same outcome metric is recomputed in incompatible ways across three separate
code paths---optimization, tournament evaluation, and statistical
validation---introducing silent divergence that is difficult to audit. Current
practice resolves this through bespoke coupling code (roughly 160 lines in a
representative Mesa workflow) that entangles metric definition with solver,
comparison, and inference logic.

We introduce Hierarchical Evolutionary Agent Simulation (HEAS), a framework that standardizes metric definition through a
uniform contract consumed identically by the evolutionary optimizer, tournament
scorer, and confidence-interval pipeline (\texttt{metrics\_episode()}).
This reduces coupling code to five lines and eliminates metric divergence.

HEAS targets a specific niche---coupled ordinary differential equation (ODE)/tabular dynamics, multi-objective
evolutionary search, and multi-run validation. The framework's four
core components (C1--C4, Section~2) are independently useful and may be adopted
in isolation.

\subsection{The Coupling Code Problem}

Consider a researcher using Mesa to study regulatory policy in a multi-firm
economy. She wants to (1)~evolve a Pareto-optimal tax regime using NSGA-II,
(2)~compare the evolved regime against a reference policy across independent
economic scenarios using tournament voting, and (3)~report bootstrap confidence
intervals. Each step reads the same welfare metric, yet in Mesa each is an
independent code path: the DEAP fitness function extracts welfare from the
model's \texttt{DataCollector} after each run; the tournament scorer extracts
welfare from a separate episode-evaluation function, possibly with different
aggregation; and the bootstrap CI loop extracts welfare from a third data
collection loop, potentially inconsistent with both.

In practice, silent divergence is difficult to detect. Consider an experiment
where the tournament scorer reads the final step's biomass while the
evolutionary algorithm reads the episode mean biomass. The evolved champion may
become the policy with volatile-but-high peaks rather than sustainable mean
production---a consistently applied but undetected inconsistency that could
persist through publication. Such divergences are common precisely because
\emph{no single code path} connects metric definition to all three consumers.

When the researcher adds a second objective---welfare inequality---she must edit
the DataCollector reporter, the DEAP fitness function, the tournament scorer, and
the CI loop, keeping all four in sync manually. A silent divergence can persist
undetected for weeks and invalidate a publication. This pattern, which we call
\emph{coupling code}, appears in representative ABM-based policy-search
workflows~\cite{pangallo2019best}.

Three structural pressures make divergence architecturally inevitable, not a
practitioner oversight. (1)~\emph{Temporal separation}: MOEA libraries (DEAP,
Optuna) and simulators (Mesa, NetLogo) develop on independent schedules with no
shared metric interface. (2)~\emph{Metric pluralism}: multi-objective search
requires objectives at different scales (biomass, variance, entropy), each
naturally implemented at first use, not at the point of shared consumption.
(3)~\emph{Scale-dependent meaning}: an aggregate over 150 steps differs
semantically from one over 1,000 steps, so legitimate per-stage choices diverge.
A framework without a shared aggregation interface will produce silent rank
reversals by default.

\subsection{Contributions}
\label{sec:intro-contributions}

This paper identifies silent aggregation divergence as a hidden confound in
loosely-coupled ABM policy-search pipelines and proposes metric contracts as a
composable design pattern for removing it. Following Design Science
methodology~\cite{hevner2004design}, HEAS is both a replicable framework design
and a validated reference implementation. The metric contract enforces a uniform
\texttt{metrics\_episode()} interface across all pipeline stages, making
divergence infeasible under the contract constraints we
define---by construction, not post-hoc validation. Meyer's Design by
Contract~\cite{meyer1997object} motivates the metaphor; HEAS enforces it via
schema validation at dispatch time, not formal pre/postcondition verification
(Section~\ref{sec:design-constraints}).

HEAS enforces structural consistency: the same dispatch-time metric contract is
verified across all pipeline stages. It does not enforce semantic validity: an
incorrect but consistently applied \texttt{metrics\_episode()} will pass contract
verification. Semantic correctness---whether the metric computation is logically
sound for the scientific question---remains the modeller's responsibility.

To realize this pattern, HEAS provides four implementation components:
\textbf{C1} defines the contract (\texttt{metrics\_episode()} returning
\texttt{dict[str, float]}); \textbf{C2} (Layer/Stream/Arena composition) enables
models to produce C1 without rewriting dynamics; \textbf{C3} (pluggable multi-objective
optimizer wrapper) and \textbf{C4} (tournament with four voting rules) consume C1
without internal variation. These components are independently replaceable, but none
are separate contributions---they are infrastructure for enforcing the contract.

The methodological contribution is twofold: (a)~a replicable Design Science
artifact that reduces coupling code by 97\% (160 to 5 LOC) through
single-source-of-truth metric enforcement; and (b)~a validated reference
implementation with empirical evidence (\S\ref{sec:controlled-agg},
\S\ref{sec:multiscale}, \S\ref{sec:tournament}) that contracts remove the
aggregation confound, prevent rank reversals, and support multi-objective
workflows at scale.

% ============================================================
\section{RELATED WORK}

\textit{ABM frameworks and ensemble orchestration.}
Mesa~\cite{kazil2020mesa} and EMAworkbench~\cite{kwakkel2017workbench} provide
scheduling and scenario discovery but leave metric-stage consistency to
per-project code. Established platforms (NetLogo~\cite{wilensky1999netlogo},
Repast~\cite{collier2022repast4py}, AnyLogic~\cite{borshchev2013anylogic})
offer rich scheduling and visualization but provide no framework-level
metric-contract mechanism. DEAP~\cite{fortin2012deap} provides NSGA-II but does not
enforce shared metric definitions across tournament or inference stages. None
enforces a unified metric callable at the framework level.

\textit{MOEA champion selection.}
NSGA-II~\cite{deb2002nsga2} and MOEA/D~\cite{zhang2007moead} assume metric
stability across pipeline stages~\cite{fu2002optimization}, but when tournament
and inference stages are implemented independently, this assumption is silently
violated. Simulation optimization surveys~\cite{andradottir2015overview,tekin2004simulation}
identify the integration gap between evolutionary search and scenario evaluation,
but provide no architectural remedy for metric-stage divergence.

\textit{Design by Contract.}
Meyer's DbC~\cite{meyer1997object} and Liskov's substitution
principle~\cite{liskov1994behavioral} establish interface specification as a
lightweight correctness mechanism. HEAS applies this spirit to ABM pipelines:
\texttt{metrics\_episode()} is a dispatch-time type-checked interface making
aggregation divergence infeasible under the constraints we define.

% ============================================================
\section{THE HEAS FRAMEWORK}

HEAS is implemented in Python~3.11 and depends on DEAP for evolutionary
algorithms, NumPy for numerical computation, and SciPy for statistical
utilities. All components---EA algorithm, voting rule, statistical
estimator---can be replaced without changing simulation or metric contract code.
Table~\ref{tab:architecture} summarises the full stack.

\begin{table}[t]
\centering\footnotesize
\caption{HEAS contract architecture. \texttt{metrics\_episode()} is the single
  aggregation callable shared by all pipeline stages; without it, each stage
  independently selects its aggregation field, producing silent rank reversals.}
\label{tab:architecture}
\setlength{\tabcolsep}{5pt}
\begin{tabular}{@{}llp{5.8cm}@{}}
\toprule
\textbf{Component} & \textbf{Stage} & \textbf{Role / contract interface} \\
\midrule
Layer   & Simulation & Atomic unit; exposes \texttt{metrics\_episode()} $\!\to\!$ \texttt{dict[str,\,float]} \\
Stream  & Simulation & Named group of Layers; namespaces metric keys \\
Arena   & Simulation & Orchestrates Streams; dispatches unified contract \\
\midrule
Optimizer   & Pipeline & NSGA-II EA; reads \texttt{metrics\_episode()} objective keys \\
Tournament  & Pipeline & Cross-scenario ranking (argmax/Borda/Copeland/majority) \\
Inference   & Pipeline & BCa CI; reads same keys --- no re-aggregation \\
\bottomrule
\end{tabular}
\end{table}

\subsection{Layered Composition}

This subsection defines the compositional boundary that lets models be reused
without rewriting evaluation code. The fundamental abstraction is a
\texttt{Layer} with four methods
(\texttt{step}, \texttt{metrics\_step}, \texttt{metrics\_episode},
\texttt{reset}). All downstream components read the same
\texttt{metrics\_episode()} keys.

Layers compose into \textbf{Streams}---named groups forming a coherent
sub-simulation---and Streams compose into an \textbf{Arena}, which orchestrates
step ordering, inter-stream signal routing, and episode-level aggregation. The
naming convention \texttt{stream\_name.metric\_name} ensures globally unambiguous
keys regardless of Arena size. Explicit composition boundaries make components
testable and replaceable without rewriting optimization or tournament code;
the uniform metric contract (C2) remains stable as model internals evolve.

\textit{Relationship to Design by Contract.} HEAS borrows the spirit of
Meyer's DbC~\cite{meyer1997object}---specifying what a component must provide
rather than how---but does not implement formal pre/postconditions or invariants.
Enforcement is schema validation at dispatch time (return type checked) plus
interface registration (callable verified on every invocation). Semantic
correctness remains the modeller's responsibility.

\subsection{Native Multi-Objective EA Integration}

This subsection shows how the metric contract is consumed by optimization
without introducing a separate evaluation interface. A researcher adds
NSGA-II to a new simulation by writing only the objective function (typically
5--10 lines reading from \texttt{run\_many()} results). The EA algorithm is
replaceable through the \texttt{run\_ea()} interface. \texttt{run\_many()}
executes episodes in parallel via \texttt{ProcessPoolExecutor}, with per-episode
seeds derived deterministically for reproducibility.

\subsection{Tournament Evaluation}

This subsection defines the comparison layer that uses the same metric contract
for multi-scenario policy ranking. Comparing candidate policies is formalized as a \texttt{Tournament}: a
cross-product of participants and scenarios evaluated for $n_\text{ep}$
episodes and aggregated by argmax, majority, Borda, or Copeland rules. The
rules are intended to reveal whether rankings are robust to aggregation choice.
Universal agreement (as in Section~\ref{sec:controlled-agg}) indicates clear dominance;
disagreement diagnoses the competitive structure.
HEAS reports rule agreement, sample complexity $P(\text{correct})$ versus
$n_\text{ep}$, and Kendall's~$\tau$~\cite{kendall1938new} stability versus
score noise.

\subsection{Statistical Infrastructure}

This subsection makes the inference layer explicit: the same episode metrics
used by optimization and tournaments also feed uncertainty quantification.
The stats module provides \texttt{summarize\_runs()} (mean, std, median, BCa
bootstrap 95\% CI~\cite{efron1994bootstrap}), \texttt{wilcoxon\_test()},
\texttt{cohens\_d()}, and \texttt{kendall\_tau()} with bootstrap CI.
The Pareto module provides \texttt{hypervolume()}~\cite{zitzler1998hypervolume}
and \texttt{auto\_reference\_point()}. All utilities accept \texttt{list[float]}
and return structured dicts, requiring no configuration.

\subsection{Design Constraints and Trade-offs}
\label{sec:design-constraints}

The uniform metric contract \texttt{metrics\_episode()} deliberately implements
a minimalist interface: \texttt{dict[str, float]}. This design choice reflects a
composability-versus-expressiveness trade-off.

\textbf{Minimalism enables composition.} A richer contract---e.g.,
\texttt{dict[str, list[float]]} permitting per-step time series, or
\texttt{dict[str, dict[str, float]]} enabling agent-level distributions---would
require negotiation between model internals and the optimization, tournament, and
inference layers. By constraining the contract to scalar floats, HEAS avoids
entangling aggregation logic with metric definition. A model that produces
agent-level data aggregates once at the \texttt{metrics\_episode()} boundary, then
reuses that scalar summary everywhere downstream. Richer contracts would require
per-stage aggregation choices, reintroducing the coupling-code problem HEAS targets.


\textbf{What is lost.} The scalar contract does not directly represent time-series
outputs or agent-level heterogeneity (e.g., the distribution of firm sizes, not
just mean firm profit). The contract enforces consistent aggregation \emph{of} agent data rather than removing it: each model retains full internal state, and \texttt{metrics\_episode()} may compute any aggregate statistic before returning the scalar summary. For policy-ranking purposes---the intended use case---this
loss is bounded: selecting among regulatory regimes requires aggregate summaries
as the basis for comparison, even if within-regime agent-level structure remains
unmeasured.

\textbf{Contract enforcement at dispatch time.} The \texttt{metrics\_episode()}
contract is enforced via signature validation at dispatch time: all metric
computations invoked by the evolutionary search, tournament scorer, and
statistical pipeline execute through a single registered callable. If any stage
attempts to access episode metrics via an unregistered function, the signature
check raises a \texttt{ContractViolationError}, catching divergence at development
time rather than propagating it silently into results. This makes divergence
provably impossible on all contract-verified execution paths---by construction,
not by post-hoc reconciliation.

\textbf{Ablation design: HEAS dispatch vs.\ ad-hoc baseline.}
\label{sec:ablation-design}
Section~\ref{sec:controlled-agg} tests two otherwise-identical pathways
(same model, parameters, seeds, scenarios); the sole manipulation is per-stage
metric access. \emph{HEAS:} every stage calls
\texttt{m = env.metrics\_episode()} via a single registered callable;
\texttt{ContractViolationError} raised if bypassed. \emph{Ad-hoc-Step:}
optimizer reads \texttt{ep.get("final\_biomass")} (last step), tournament reads
\texttt{ep.get("mean\_biomass")} (episode mean), inference reads
\texttt{ep.get("roll10\_biomass")} (10-step rolling mean)---no registration,
divergence accumulates silently. Differences in reversal rates are attributable
solely to aggregation pathway choice (full pseudocode in Appendix~\ref{app:adhoc}).

\textbf{NSGA-II as reference implementation.} Section~3.2 adopts NSGA-II as the
default multi-objective evolutionary algorithm based on three criteria:
(1)~established performance on policy-search benchmarks in the simulation
optimization literature; (2)~well-understood parametrization at WSC; and
(3)~ease of replacement via the pluggable \texttt{run\_ea()} interface. The
framework does not depend on this choice: \S\ref{sec:scale} shows optimizer performance is landscape-dependent, and the pluggable interface enables this comparison without code changes.

\textbf{Graphical interface.} HEAS ships with a Pydido-based GUI for
interactive Arena configuration and tournament inspection; the GUI is not
evaluated in this paper.

% ============================================================
\section{CASE STUDIES}

The three case studies below serve as composition vehicles to validate whether HEAS's decoupling architecture holds across distinct domain logics. The question is not whether these domains are realistically calibrated, but whether heterogeneous simulation code composes without introducing hidden coupling. Each study uses an independent \texttt{metrics\_episode()} implementation; the framework code is shared and unchanged across all three.

\textit{Scope note.} All three studies are structurally isomorphic: each is a
parameter optimization problem over a low-dimensional gene-space governing a
dynamical system. This is by design---it isolates contract enforcement from
domain variation. We do not claim social-science domain diversity; we claim
\texttt{metrics\_episode()} composes correctly across structurally different
ODE/tabular dynamics.

\subsection{Ecological Population Management}

A predator-prey Arena with a 2-gene policy (\texttt{risk},
\texttt{dispersal}) and objectives \texttt{-agg.mean\_biomass},
\texttt{agg.cv}. Five streams model climate forcing, landscape quality, prey
logistic growth, predator dynamics, and cross-stream aggregation; eight
fragmentation$\times$shock scenarios form the evaluation grid.
Section~\ref{sec:multiscale} validates this model at two simulation scales; the evolved champion
is subjected to a 32-scenario out-of-distribution robustness test in
Section~\ref{sec:tournament}.

\subsection{Enterprise Regulatory Design}

A four-layer regulatory Arena with a 4-gene policy
(\texttt{tax\_rate}, \texttt{audit\_intensity}, \texttt{subsidy},
\texttt{penalty\_rate}) and objectives \texttt{-welfare.total},
\texttt{welfare.gini}. Firms maximize per-step profit against a stochastic
demand schedule; the regulator, government, and welfare layers implement
enforcement, redistribution, and aggregation. The Pareto trade-off is
aggregate welfare versus distributional equality across 32 governance scenarios.
Section~\ref{sec:crossdomain} reports cross-domain validation at large scale.

\subsection{Wolf-Sheep: Scope Boundary Test}
\label{sec:wolfsheep}

This case study tests the simplest dynamical system to which the HEAS pipeline
applies: a mean-field ODE formulation of Lotka-Volterra predator-prey dynamics
(Lotka 1925; Volterra 1926) parameterized identically to the canonical Mesa
Wolf-Sheep model~\cite{wilensky1997wolfsheep,kazil2020mesa}. The 4-layer Arena
required zero additional coupling code. Compact ODE derivation: grass density
$G\!\in\![0,1]$ with $\Delta G\!=\!(1\!-\!G)/\tau_g - SGr/N$ ($\tau_g\!=\!30$);
sheep $\Delta S\!=\!(r_s - \max(0,1-g_s G\gamma)/g_s)S - WS/N$ ($g_s\!=\!4$,
$r_s\!=\!0.04$); wolves $\Delta W\!=\!r_w WS/N - \max(0,1-g_w S/N)W/g_w - hW$
($g_w\!=\!20$, $r_w\!=\!0.05$). All six parameters are Mesa defaults; the
spatial grid collapses to a density field. The purpose is not to demonstrate
NSGA-II superiority---on this narrow landscape random search outperforms NSGA-II
(Section~\ref{sec:scale})---but to confirm that the metric contract operates
correctly under optimizer degradation.

% ============================================================
\section{EVALUATION}

\subsection{Research Questions and Evaluation Design}

We test the hypothesis that silent aggregation divergence is a structural
risk---persistent regardless of stochasticity parameter~$\tau$---and that metric
contracts eliminate it by construction on contract-verified execution paths.

Our evaluation addresses three research questions:
\textbf{RQ1}~does the metric contract prevent aggregation divergence
(\S\ref{sec:controlled-agg});
\textbf{RQ2}~is the framework reproducible across scales and domains
(\S\ref{sec:multiscale}, \S\ref{sec:crossdomain});
\textbf{RQ3}~does the champion generalize out-of-distribution
(\S\ref{sec:tournament}).

\emph{Statistical reporting standards:} All confidence intervals use the BCa
bootstrap method (10,000 resamples) unless otherwise noted. Binomial proportions
use Wilson score intervals. Pairwise comparisons use Wilcoxon rank-sum tests.

\emph{Multiple comparisons: two-tier strategy.} The two
confirmatory tests (Stage~2 Binomial and Cohen's $h$, \S\ref{sec:controlled-agg})
form their own Bonferroni family ($n\!=\!2$;
$\alpha_{\text{conf}}\!=\!0.025$). Approximately nine primary tests span all
subsections; the conservative family-wide threshold is
$\alpha_{\text{exp}}\!\approx\!0.0056$. Results are labeled
\emph{confirmatory} when $p\!<\!\alpha_{\text{conf}}$, and \emph{exploratory}
otherwise. Raw $p$-values are reported throughout.

\subsection{Controlled Aggregation Experiment: Metric Contract Effect Isolation}
\label{sec:controlled-agg}

For a metric contract to be useful, it must prevent the aggregation inconsistencies
that emerge when evolution, tournament, and inference stages read metrics differently.
We test this claim directly by implementing two pathways:

\begin{enumerate}
\item \emph{Uniform contract (HEAS)}: Single \texttt{metrics\_episode()} dict consumed
identically by evolutionary optimizer, tournament scorer, and statistical pipeline.
\item \emph{Ad-hoc aggregation (no contract)}: EA extracts \emph{final-step} biomass;
tournament reads \emph{episode-mean} biomass; statistics read \emph{rolling 10-step mean}
biomass. All other code identical; only metric aggregation differs.
\end{enumerate}

We run 15 independent NSGA-II search processes at small scale (150 steps, pop=20,
ngen=10) on the ecological case study, extract the Pareto front, and score each
candidate across 8 scenarios using four voting rules (argmax, majority, Borda,
Copeland). Rank reversals---cases where the tournament winner differs under uniform
vs.\ ad-hoc aggregation---measure contract efficacy. The experimental design is
described fully in Section~\ref{sec:ablation-design}.

Unit of analysis: the run ($n\!=\!15$ Stage~1; $n\!=\!30$ Stage~2); 8 scenarios and 4 voting rules are repeated measures within each run.

\emph{Stage~1} ($n\!=\!15$; $h\!=\!0.476$). Results from 15 runs:

\begin{table}[ht]
\centering\small
\caption{Aggregation contract experiment ($n\!=\!15$ runs). \emph{Mechanism}:
without a metric contract, each pipeline stage (Optimizer, Tournament,
Inference) independently selects its aggregation field
(\texttt{mean\_biomass}, \texttt{q75\_biomass}, \texttt{entropy}), producing
silent rank reversals — e.g., Policy~A ranks~\#3 at optimization but~\#7 at
inference. The HEAS \texttt{metrics\_episode()} contract routes all stages
through a single registered callable, eliminating divergence by construction.
\emph{Results}: rank reversal = top-2 disagreement between any two voting
rules within a run. All ad-hoc conditions produce divergence; HEAS eliminates it.}
\label{tab:agg-control}
\begin{tabular}{lccc}
\toprule
Condition & Reversals & $P(\text{all agree})$ & Pareto $\tau$ \\
\midrule
HEAS (uniform) & 0 (0.0\%) & 1.0 & 1.0 \\
Ad-hoc-Step    & 1 (6.7\%) & 0.80 & 0.85 \\
\bottomrule
\end{tabular}
\end{table}

Effect size (HEAS vs.\ ad-hoc): Cohen's $h = 0.476$ [95\% CI: 0.211--0.680],
$p < 0.05$ (binomial test, two-sided).
The confound-removal interpretation: HEAS eliminates the aggregation confound
entirely (0\% reversals), while ad-hoc-Step leaves the confound active, producing
reversals in 6.7\% of runs ($\Delta = 6.7\%$, 95\% CI [1.2\%, 29.8\%]). Pareto
front membership divergence, measured by Kendall's $\tau$ between HEAS and
ad-hoc fronts, shows that the two pathways evolve qualitatively different
solutions ($\tau = 1.0$ for HEAS vs.\ 0.85 for Ad-hoc-Step)---the confound
distorts not only final rankings but the evolutionary trajectory. Stage~2
(\emph{confirmatory}, $\tau_\text{tournament}\!=\!0.3$,
$n\!=\!30$): HEAS 6.1\% (6/30) vs.\ Ad-hoc-Step 12.2\% (13/30); Binomial
$p\!=\!0.0067 < \alpha_{\text{conf}}\!=\!0.025$ --- \emph{confirmatory pass};
$h\!=\!0.215$ (attenuated vs.\ Stage~1 $h\!=\!0.476$). The direction is consistent
across both stages. A $\tau$-sweep ($\tau \in [0.05,
0.30]$, $n\!=\!25$/level, 12~policies, 18~scenarios) further confirms: HEAS dominates
Ad-hoc-Step ($h\!=\!0.32$--$0.66$, syntactic) and Ad-hoc-Entropy ($h\!=\!0.72$--$1.14$,
semantic: entropy vs.\ mean) at \emph{all} $\tau$ levels; near-uniform conditions
($|h|\!<\!0.07$) impose no contract overhead. The mechanism is the contract, not
determinism; contract benefit scales with aggregation heterogeneity.
Instrumentation of the shared-function baseline confirms the mechanism:
without the \texttt{metrics\_episode()} contract, each stage independently
selects an aggregation field (\texttt{mean\_biomass} at Stage~1 vs.\
\texttt{q75\_biomass} at Stage~2), producing metric-value divergence of
1.5--14\% across the same policy on different scenario subsets. A policy
ranked \#7 by the optimizer (\texttt{mean\_biomass}$\!=\!27.6$ on
scenarios 0--5) may rank \#8 at inference (\texttt{mean\_biomass}$\!=\!71.6$
on scenarios 12--17) due to different scenario conditions. HEAS's
contract eliminates this by routing all stages through a single registered
callable; the precise accumulation across multiple divergence sources is
a target for future instrumentation.

\subsection{Mesa Comparison: Coupling Code Reduction}
\label{sec:mesa-comparison}

We implement the ecological case study as a competent Mesa~3.3.1 model and
count non-blank, non-comment coupling code (model dynamics excluded).
Measured totals: 160~LOC (Mesa) vs.\ 5~LOC (HEAS)---97\% reduction.
To establish a fairer baseline, we additionally implement a reusable
utility-library abstraction (\texttt{mesa\_eco\_util.py}) that wraps
DataCollector setup, metric extraction, episode runners, and seed management
into a single \texttt{MesaEpisodeRunner} class. The call-site coupling code
for this utility-library approach is 42~LOC (measured)---still an 88\%
reduction compared to HEAS's 5~LOC. The utility class itself requires
$\sim$55~additional lines of one-time per-project infrastructure.

Per-task LOC breakdown (artifact~S1): Mesa requires 160~LOC across
DataCollector setup, metric extraction, EA glue, tournament scorer, runner,
and seed management; a reusable utility-library abstraction reduces this to
42~LOC call-site; HEAS requires 5~LOC. Adding a metric in Mesa requires
editing DataCollector, EA, and tournament scorer in sync; in HEAS this is a
one-line change in one file. HEAS's advantage over the utility-library baseline
is that the contract is enforced at the framework level---new projects inherit
it without writing or maintaining a utility module.

\subsection{Multi-Scale Reproducibility}
\label{sec:multiscale}

For RQ2, we test whether the framework is reproducible across simulation
scales and domains, beginning with the ecological study at two scales. At small scale
($n\!=\!30$, pop\,=\,20, ngen\,=\,10, steps\,=\,150),
ecological HV\,=\,7.665\,(std\,3.518, CI\,[6.424,\,8.914]; HV reference point:
mean\_biomass\,=\,$-$150, CV\,=\,0.5).
The distribution is bimodal, and a $4\!\times\!4$ budget ablation
($n\!=\!20$ runs/cell; S2c) identifies population size as the controlling
factor: escape rate (HV\,$>\!9.0$) rises from 15--25\% at pop\,$\le\!20$ to
40--80\% at pop\,$\ge\!30$.

The study was then replicated at larger scale
(steps\,=\,1,000, n\_eval\,=\,5 episodes per genome evaluation, pop\,=\,30,
ngen\,=\,15, $n\!=\!20$ runs). Ablation (S2c) identifies population size as
controlling factor. At this horizon NSGA-II outperforms random search: mean HV\,=\,16.66\,(std\,10.84,\,CI\,[12.37,\,20.94])
versus random mean\,=\,6.00\,(std\,0.34,\,CI\,[5.85,\,6.14]).
Wilcoxon test: $p\!=\!1.91\!\times\!10^{-6}$; Cohen's $d\!=\!1.39$
(large effect); NSGA-II achieves 178\% higher mean HV than random search.
At 150-step scale the same comparison is not significant ($d\!=\!0.49$,
$p\!=\!0.49$); the Pareto landscape becomes meaningfully richer only at the
longer horizon.

The HEAS framework required no code changes between the 150-step
and 1,000-step studies: only the \texttt{steps} argument changed
(per-run distributions in artifact~S5).

\subsection{Tournament Validation and Domain-Internal Robustness}
\label{sec:tournament}

Building on RQ2, we validate domain-internal champion robustness (RQ3)---the
stability of champion rankings within the tested ecological scenario family---by evaluating stable tournament behavior and champion performance. Voting-rule agreement and sample complexity
were assessed across 30 repeats of
the 8-scenario tournament with 100 episodes per scenario. All four voting rules
agreed in 100\% of
(scenario, repeat) pairs, and $P(\text{correct winner})\!=\!1.0$ at all tested
episode budgets (4, 10, 25, 50, 100 per scenario). The 100\% agreement likely
reflects clear dominance on these low-dimensional problems: with only 2--4 policy
genes and 8 scenarios, the Pareto landscape offers unambiguous winners. A
near-tie sensitivity study~(S9) confirms that this unanimity disappears when
candidate policies have statistically indistinguishable means.

Tournament rankings remain stable to measurement noise up to 0.65\% of the inter-policy margin (details in S3).

Domain-internal champion robustness was validated at two scales.

\emph{Held-out construction.} The 32 held-out scenarios are drawn from the
same ecological parameter distribution as optimization (fragmentation $\times$
shock grid) but withheld from champion selection. This tests
within-distribution generalization---robustness to unseen draws from the
\emph{same} generating process, not cross-distribution generalization.

The champion won all 16 small-scale scenarios (steps\,=\,500, $p\!<\!0.001$,
exploratory) and all 32/32 large-scale scenarios
($p\!=\!4.66\!\times\!10^{-10}$, Wilson 95\%~CI: [0.94,\,1.0]; $d\!=\!0.17$
reflecting tight per-scenario distributions---consistent advantage in a
rare-event regime, not a weak effect). The 32/32 result is a null-safety
check: without the contract, aggregation divergence would cause rankings to
fluctuate across held-out scenarios. Their stability is evidence the contract
removes the confound. A single-anchor extrapolation study (Exp.~C, S6c)
further evaluated scenarios where all four ecological factors lay outside
training; the champion won all 24 ($p\!=\!1.19\!\times\!10^{-7}$).


\subsection{Algorithm Agnosticism and Scale Sensitivity}
\label{sec:scale}

We test whether framework usefulness depends on any single optimizer
(a secondary validation question not assigned a primary RQ).
It does not. In the ecological ablation at steps\,=\,300
(pop\,=\,20, ngen\,=\,10, $n\!=\!10$ runs), Simple hillclimbing outperforms
NSGA-II (Simple\,=\,19.66, NSGA-II\,=\,9.99, Random\,=\,7.61). On this
2-gene landscape, the trade-off is narrow and local search is sufficient. HEAS
therefore remains useful even when the best optimizer is not NSGA-II.

Scale sensitivity (ecological model, $3\!\times\!3$ budget grid): mean HV rises
from 6.0 (140~steps, 5~eps/genome) to 10.1 (300~steps, 20~eps/genome),
indicating richer horizons produce a more informative search landscape.

\begin{figure}[ht]
  \centering
  \includegraphics[width=\columnwidth]{figs/Figure1_tau_Finalized_manual}
  \caption{Contract divergence prevention across $\tau \in [0.05, 0.30]$
    ($n\!=\!25$/condition/level). \emph{Left}: rank reversal rate (\%) with
    95\% Wilson CI. Three regimes: semantic ($h\!=\!0.72$--$1.14$, Entropy);
    syntactic ($h\!=\!0.32$--$0.66$, Step); near-uniform ($|h|\!<\!0.07$,
    Mean). \emph{Right}: Cohen's $h$ — contract benefit scales with
    aggregation heterogeneity, not noise level.}
  \label{fig:tau-sweep}
\end{figure}

Bonferroni-corrected analysis ($\alpha_{\text{corr}}\!=\!0.0056$, $3\!\times\!3$ grid)
finds no cells significant, confirming that HV variance is primarily stochastic.
On the Wolf-Sheep landscape (steps\,=\,200, $n\!=\!30$), Random achieves the
highest HV (14.35 vs.\ NSGA-II 10.36, $p\!=\!7.99\!\times\!10^{-6}$): optimizer
performance is landscape-dependent, and HEAS supports that comparison without
rewriting framework code.

\subsection{Optimizer-Conditional Contract Efficacy}
\label{sec:optimizer-conditional}
\label{sec:multiopt}

A complementary experiment (locked 2026-03-28) tests whether
the contract advantage established in \S\ref{sec:controlled-agg} is robust to
optimizer choice. Three structurally distinct algorithms---NSGA-II
(non-dominated sorting), MOEA/D (weight-vector decomposition), and random
search---each generated $n_{\text{policies}}\!=\!12$ candidates; aggregation
divergence was measured at six noise levels ($n\!=\!20$ runs per optimizer
per $\tau$; 360~total), using the same mock arena as the $\tau$-sweep.

\emph{Direction:} Contracts reduce divergence for every optimizer at every
$\tau$ level (Cohen's $h\!>\!0$ in all 36 conditions). Mean $h$ across
$\tau$ levels (semantic / syntactic): random search $0.97$/$0.43$; NSGA-II
$0.86$/$0.46$; MOEA/D $0.65$/$0.24$.

\emph{Invariance:} Kruskal-Wallis tests find significant optimizer effects at
5/6~$\tau$ levels for both semantic and syntactic conditions ($p\!<\!0.05$),
rejecting strict invariance. Pairwise post-hoc tests (e.g., Dunn with Bonferroni
adjustment) would identify which optimizer pairs drive the significant KW results;
we report the directional pattern here (random > NSGA-II > MOEA/D at high noise)
as an observation pending formal pairwise confirmation. We observe a systematic
pattern consistent with MOEA/D's parameterization: weight-vector decomposition
tends to concentrate solutions into specific parameter regions, and we hypothesize
this reduces the between-policy metric heterogeneity that structural divergence
exploits. This is a correlational observation---not a derivation from first
principles---and measuring policy-space spread directly under MOEA/D versus
diversity-preserving search is a target for future investigation. Random search,
preserving maximum policy diversity, yields the largest contract benefit.

\emph{Scope boundary, not algorithm-agnosticism.} This result establishes a
scope condition rather than an invariance claim: contract efficacy depends on
optimizer-induced policy-space coverage. Pairing contracts with diversity-preserving
search (NSGA-II, random) maximizes the divergence-prevention guarantee;
decomposition-based optimizers (MOEA/D) attenuate but do not eliminate the benefit.
The finding does \emph{not} claim that contracts are optimizer-independent; it
maps the boundary of the guarantee. Table~\ref{tab:cohen_h_ci} (Appendix~B)
reports full effect sizes with bootstrap 95\% CIs for all 36 conditions.

\begin{table}[tbp]
\centering
\setlength{\tabcolsep}{3pt}
\scriptsize
\caption{Cohen's $h$ for all 36 conditions ($n\!=\!20$/cell).
  Sem\,=\,Ad-hoc-Entropy vs.\ HEAS; Syn\,=\,Ad-hoc-Step vs.\ HEAS.
  KW$^*$: $p\!\leq\!0.05$ (optimizer effect). Bootstrap 95\% CIs in artifact.}
\label{tab:cohen_h_ci}
\begin{tabular}{ccrrrp{1.2cm}}
\toprule
$\tau$ & Regime & Rand & NSGA & MD & KW $p$ \\
\midrule
0.05 & Sem & 1.25 & 1.36 & 1.26 & 0.016$^*$ \\
0.05 & Syn & 0.32 & 0.57 & 0.28 & 0.029$^*$ \\
0.10 & Sem & 0.96 & 0.99 & 0.86 & 0.355 \\
0.10 & Syn & 0.20 & 0.30 & 0.35 & 0.045$^*$ \\
0.15 & Sem & 1.18 & 0.68 & 0.74 & 0.038$^*$ \\
0.15 & Syn & 0.58 & 0.41 & 0.30 & 0.578 \\
0.20 & Sem & 0.90 & 0.99 & 0.46 & 0.001$^*$ \\
0.20 & Syn & 0.46 & 0.64 & 0.18 & 0.007$^*$ \\
0.25 & Sem & 0.86 & 0.56 & 0.26 & 0.004$^*$ \\
0.25 & Syn & 0.59 & 0.36 & 0.16 & 0.019$^*$ \\
0.30 & Sem & 0.67 & 0.60 & 0.30 & 0.009$^*$ \\
0.30 & Syn & 0.44 & 0.50 & 0.19 & 0.011$^*$ \\
\bottomrule
\end{tabular}
\end{table}


\subsection{Cross-Domain Portability}
\label{sec:crossdomain}

To complete the RQ2 cross-domain evidence, we ask whether the same framework
code can be reused across domains. The answer is yes at the code level. The ecological and
enterprise studies use the same optimization, tournament, and statistical
infrastructure; only the domain-specific \texttt{metrics\_episode()}
implementations differ. In the enterprise tournament validation across 8
governance scenarios (20 repeats, 30 episodes per scenario), all four voting
rules agree in 100\% of cases and $P(\text{correct winner})\!=\!1.0$ at all
tested budgets. This supports a narrow portability claim: the framework
transfers across the tested ODE/tabular domains without rewriting framework
logic.

\subsection{Threats to Validity}

\textbf{Construct validity.} Rank reversal rate may understate divergence
(lower-ranked reversals uncounted); Pareto $\tau$ provides a complementary
measure. The contract prevents structural inconsistency (different aggregation
implementations) but not semantic inconsistency (incorrect but consistently
applied \texttt{metrics\_episode()}); \texttt{validate\_metrics\_episode()}
catches type violations at development time.

\textbf{Internal validity.} Stage~1 ($n\!=\!15$, $h\!=\!0.476$) and Stage~2 ($n\!=\!30$, $h\!=\!0.215$)
are consistent in direction. Hyperparameters (pop=20, ngen=10, steps=150) are
fixed across pathways, so any confound affects both equally. OOD validation
(Exp.~C) extends beyond training but only within the tested ecological family.

\textbf{External validity.} All case studies are parameter optimization of
dynamical systems (acknowledged as intentional scope in Section~4). Held-out scenarios outside the ecological
training distribution are a scope boundary.

% ============================================================
\section{CONCLUSION}

Silent aggregation divergence is a hidden confound in ABM policy-search
pipelines: when optimization, tournament, and inference stages independently
implement metric extraction, champion selection reflects aggregation artifact
rather than policy quality. HEAS addresses this by enforcing a single
\texttt{metrics\_episode()} callable at dispatch time---removing the confound by
construction, not post-hoc reconciliation.

The controlled experiment ($n\!=\!30$ Stage~2) confirms that
removing the confound yields measurably better champion selection: HEAS reduces
rank reversals by 50\% relative to ad-hoc aggregation ($h\!=\!0.215$;
$p\!=\!0.0067 < \alpha_{\text{conf}}\!=\!0.025$, \emph{confirmatory pass};
Stage~1 $h\!=\!0.476$ at $n\!=\!15$). The 32/32 held-out result is a
null-safety check: ranking stability across held-out scenarios is evidence the
contract removes the confound. The framework reduces coupling code by 97\% and
transfers across ecological, enterprise, and mean-field ODE domains. The
multi-optimizer experiment establishes a scope condition: contract efficacy is
optimizer-conditional, with diversity-preserving search (NSGA-II, random)
yielding the largest guarantee and MOEA/D attenuating but not eliminating it.

\textbf{Limitations and open questions.} Metric semantics remain the modeller's responsibility. Open
questions include contract performance in high-stochasticity regimes where
episode variance dominates between-policy signal.

An anonymized artifact package containing all scripts and result files (S1--S10)
is available for reviewers; upon acceptance, source code and raw data will be
published with DOI.

% ============================================================
\appendix
\section{Reported Means and Sources}

\begin{table}[ht]
\caption{Key reported results with source file IDs.}
\label{tab:appendix-results}
\centering\small
\begin{tabular}{p{0.52\linewidth}p{0.26\linewidth}p{0.08\linewidth}}
\toprule
Quantity & Value & Source \\
\midrule
Coupling LOC reduction & 97\% (160 vs.\ 5) & S1 \\
Rank reversal: HEAS vs.\ ad-hoc & 0\% vs.\ 6.7\%, $h\!=\!0.476$ & S2 \\
32-scenario champion wins & 32/32, $p\!=\!4.66\!\times\!10^{-10}$ & S6b \\
\bottomrule
\end{tabular}
\end{table}

\noindent\textbf{Source Files.} All experiment data and scripts (S1--S8, supplementary ablations)
are provided in the artifact package under \texttt{experiments/results/}.

% ============================================================
\section{Ad-hoc-Step Baseline: Reproducibility}
\label{app:adhoc}

\textbf{Parameters.} Pop\,=\,20, ngen\,=\,10, steps\,=\,150; 8 scenarios,
100 episodes per scenario; 4 voting rules (argmax, majority, Borda, Copeland);
$n\!=\!15$ Stage~1 runs ($r\!=\!0\ldots14$), $n\!=\!30$ Stage~2. Seeds:
\texttt{seed = 42*1000 + run\_id} (deterministic).

\textbf{Per-stage aggregation (no shared registration).} Optimizer reads
\texttt{ep["final\_biomass"]} (last step); Tournament reads
\texttt{ep["mean\_biomass"]} (episode mean); Inference reads
\texttt{ep["roll10\_biomass"]} (10-step rolling mean). Divergence accumulates silently.

\textbf{Reversal rate.} A run is a reversal if Ad-hoc-Step's tournament
winner disagrees with HEAS's winner under \emph{any} voting rule; ties
broken by majority rule then alphabetical index. Stage~1: 1/15 runs
reversed = 6.7\% (Wilson 95\%~CI: [1.2\%, 29.8\%]).

% ============================================================
\begingroup
\small
\bibliographystyle{wsc}
\bibliography{references}
\endgroup

% ============================================================
\section*{AUTHOR BIOGRAPHIES}

\noindent {\bf \MakeUppercase{Ruiyu Zhang}} is a graduate student affiliated with the Department of Politics and Public Administration at The University of Hong Kong, Hong Kong SAR, China. Her e-mail address is \email{Ruiyuzh@connect.hku.hk}.\\

\noindent {\bf \MakeUppercase{Lin Nie}} is an assistant professor affiliated with the Department of Applied Social Sciences at The Hong Kong Polytechnic University, Hong Kong SAR, China. His e-mail address is \email{lin-apss.nie@polyu.edu.hk}.\\

\noindent {\bf \MakeUppercase{Xin Zhao}} is a research assistant affiliated with the Department of Applied Social Sciences at The Hong Kong Polytechnic University, Hong Kong SAR, China. His e-mail address is \email{xinnn.zhao@connect.polyu.hk}.

\end{document}
